\subsection{表示}

定义 1. 设 \(\mathfrak{g}\) 是 \(\mathbf{K}\) 上的李代数，\(\mathbf{M}\) 是一个 \(\mathbf{K}\) -模。从 \(\mathfrak{g}\) 到李代数 \(\mathfrak{gl}(\mathbf{M})\) 的同态称为 \(\mathfrak{g}\) 在模 \(\mathbf{M}\) 上的表示。单射表示称为忠实表示。如果 \(\mathbf{K}\) 是域，则 \(\mathbf{M}\) 在 \(\mathbf{K}\) 上的维数（有限或无限）称为该表示的维数。\(x \mapsto \operatorname{ad} x\) 在 \(\mathfrak{g}\) -模 \(\mathbf{K}\) 上的表示 \(\mathfrak{g}\) 称为 \(\mathfrak{g}\) 的伴随表示。

因此，g 在 M 上的表示是从 g 到 M 的自同态模的 K-线性映射 \(\rho\)，满足

\[
\rho ([ x, y ]). m = \rho (x) \rho (y). m - \rho (y) \rho (x). m
\]

对所有 \(x \in g, y \in g, m \in M.\)

*例. 设 G 是实李群，g 是其李代数，\(\rho\) 是 G 在有限维实向量空间 E 上的解析表示。那么，相应的从 g 到 gl(E) 的同态就是 g 在 E 上的表示。*

设 \(\mathbf{U}\) 是 \(\mathfrak{g}\) 的包络代数。第 §2 节第 1 号的命题 1 在 \(\mathfrak{g}\) 在 \(\mathbf{M}\) 上的表示与 \(\mathbf{U}\) 在 \(\mathbf{M}\) 上的表示的集合之间建立了一一对应。另一方面，我们知道（《代数》，第 VIII 章，§13，第 1 号），结合代数 \(\mathbf{U}\) 的表示这一概念与左 \(\mathbf{U}\)-模这一概念等价。

定义 2. 设 g 是 K 上的李代数，U 是其包络代数。U 上的有单位元左模称为左 g-模，或简称 g-模。

若 M 是 g-模且 \(x \in U\)，则以 \(x_{M}\) 表示由 x 在 M 上定义的自同态（参见《代数》，第 VIII 章，§ 1，第 2 号）。

U 上的有单位元右模称为右 g-模。这样的模可与左 \(U^{0}\)-模等同，即（§2，第 4 号）可与左 \(g^{0}\)-模等同。

设 \(\phi\) 是 U 的主反自同构。如果 M 是右 g-模，则在 M 上规定左 g-模结构，写作 \(a.m = m.\phi (a)\)，其中 \(m\in \mathbf{M}\)，\(a\in \mathrm{U}\)。

模论中的概念和结果可以翻译成表示的语言：

\begin{enumerate}
\def\labelenumi{(\arabic{enumi})}
\tightlist
\item
  \(\rho\) 在 \(\rho'\) 和 \(\mathfrak{g}\) 上的两个表示 \(\mathbf{M}\) 和 \(\mathbf{M}'\)，当 \(\mathfrak{g}\)-模 \(\mathbf{M}\) 和 \(\mathbf{M}'\) 同构时，称为相似的或同构的。为此，当且仅当存在一个从 K-模 \(u\) 到 K-模 \(\mathbf{M}\) 的同构 \(\mathbf{M}'\)，使得
\end{enumerate}

\[
\rho^ {\prime} (x) = u \circ \rho (x) \circ u ^ {- 1}
\]

对所有 \(x \in g\)。

\begin{enumerate}
\def\labelenumi{(\arabic{enumi})}
\setcounter{enumi}{1}
\item
  对所有 \(i \in I\)，设 \(\rho_{i}\) 是 g 在 \(M_{i}\) 上的表示。令 M 为各 g-模 \(M_{i}\) 的直和所成的 g-模。相应地有 g 在 M 上的表示 \(\rho\)，称为各 \(\rho_{i}\) 的直和，记作 \(\sum_{i\in I}\rho_{i}\)（当有 n 个表示 \(\rho_{1} + \cdots + \rho_{n}\) 时也记作 \(\rho_{1}, \ldots, \rho_{n}\)）。若 \(m = (m_{i})_{i \in I}\) 是 M 的元素且 \(x \in g\)，则 \(\rho(x).m = (\rho_{i}(x).m_{i})_{i \in I}\)。
\item
  \(\rho\) 在 \(\mathfrak{g}\) 上的表示 \(\mathbf{M}\)，当相应的 \(\mathfrak{g}\)-模是简单模时，称为简单的或不可约的。就是说，\(\mathbf{M}\) 中不存在在所有 \(\{0\}\)（\(\mathbf{M}\)）下稳定的子 K-模（除 \(\rho(x)\) 和 \(x \in \mathfrak{g}\) 外）。一类简单 \(\mathfrak{g}\)-模（《代数》，第 VIII 章，§ 3，第 2 号）定义了 \(\mathfrak{g}\) 的一类简单表示。
\item
  \(\rho\) 是 \(\mathfrak{g}\) 在 \(\mathbf{M}\) 上的表示；当相应的 \(\mathfrak{g}\)-模是半单模时，称其为半单的或完全可约的。这等价于说，\(\rho\) 与简单表示的直和相似，或者每个子 K-模
\end{enumerate}

M 在 \(\rho(x)\)（\(x \in \mathfrak{g}\)）下稳定，都有一个在 \(\rho(x)\)（\(x \in \mathfrak{g}\)）下稳定的补模（参见《代数》，第 VIII 章，§ 3，第 3 号）。

\begin{enumerate}
\def\labelenumi{(\arabic{enumi})}
\setcounter{enumi}{4}
\item
  设 \(\delta\) 是 \(\mathfrak{g}\) 的一类简单表示，对应于一类简单 \(\mathfrak{g}\)-模 C。另一方面，设 \(\rho\) 是 \(\mathfrak{g}\) 在 M 上的表示。\(\mathrm{M}_{\mathbb{C}}\)-模 M 的 C 型等型分量 \(\mathfrak{g}\)（《代数》，第 VIII 章，§ 3，第 4 号）也称为 M 的 \(\delta\) 型等型分量。该分量是 M 中所有在 \(\rho(x)\) 下稳定且 \(\rho(x)\) 在其上诱导出 \(\delta\) 类表示的子 K-模之和；它是其中某些子模的直和；若 \(\mathrm{M}_{\mathbb{C}}\) 的长度为 \(n\)，则称 \(\rho\) 含有 \(\delta\) 个 \(n\)。不同 \(\mathrm{M}_{\mathbb{C}}\) 的和是直和；当且仅当 \(\rho\) 半单时，它等于 M。
\item
  设 \(\rho, \rho'\) 是 \(\mathfrak{g}\) 的两个表示。如果 \(\rho'\) 的模是 \(\rho\) 的模的子模（分别为商模），则称 \(\rho'\) 是 \(\rho\) 的子表示（分别为商表示）。
\end{enumerate}

设 M 是 K-模。g 在 M 上的零表示在 M 上定义了 g-模结构。赋予此结构的 M 称为平凡 g-模。

设 M 是 g-模。M 的子 g-模的商 g-模也都是 M 的商模的子 g-模：它们由 M 的两个子 g-模 U、U'（满足 U ⊃ U'）以及 g-模 U/U' 得到。若上述类型的所有简单模都同构于给定的简单 g-模 N，则称 M 为 N 型纯 g-模。若 ρ 和 σ 是分别对应于 M 和 N 的 g 表示，也说 ρ 是 σ 型纯表示。

设 \(M'\) 是 M 的子 g-模。M 是 N 型纯模的充要条件是 \(M'\) 和 \(M/M'\) 都是 N 型纯模。必要性显然。设此条件成立，令 U、\(U'\) 是 M 的子 g-模，满足 \(U' \subset U\) 且 \(U/U'\) 简单；令 \(\phi\) 是从 M 到 \(M/M'\) 的典范同态；若 \(\phi(U) \neq \phi(U')\)，则 \(U/U'\) 同构于 \(\phi(U)/\phi(U')\)，因而同构于 N；若 \(\phi(U) = \phi(U')\)，则 \(U \subset U' + M'\)，所以 \(U/U'\) 同构于 \((U' + M')/U'\) 的一个简单子模，而后一个模又同构于 \(M'/(U' \cap M')\)；故 \(U/U'\) 再次同构于 N，从而 M 是 N 型纯模。

以下设 M 是 g-模，并假设 M 的 N 型纯子 g-模的集合有最大元 \(M'\)。那么，M 的每个 N 型纯子模 \(M''\) 都包含于 \(M'\) 中。因为 \(M''/(M' \cap M'')\) 和 \(M'\) 都是 N 型纯模，所以上述结果表明 \(M' + M''\) 是 N 型纯模，从而 \(M' + M'' \subset M'\)。

假设 g-模 M 有一个 Jordan-Hölder 序列 \((\mathrm{M}_{1})_{0\leqslant i\leqslant n}\)。M 是 N 型纯模的充要条件是 \(M_{0}/M_{1}\)、\(M_{1}/M_{2}\)、\(\ldots\)、\(M_{n-1}/M_{n}\) 都同构于 N；必要性显然，而充分性由上述结果对 n 作归纳立即得到。

命题 1. 设 g 是 K 上的李代数，a 是 g 的理想。设 M 是 g-模，N 是简单 a-模。把 M 看作 a-模，并假设 M 的 N 型纯子 a-模的集合有最大元 M'。则 M' 是 M 的子 g-模。

设 \(y \in \mathfrak{g}\)。令 \(\phi\) 为从 M 到 \(M / M'\) 的典范映射，令 \(f\) 为从 \(m \mapsto \phi(y_M. m)\) 到 \(M'\) 的映射 \(M / M'\)。只需证明 \(f(M') = \{0\}\)。设 \(x \in \mathfrak{a}\)。则对 \(m \in M\)，

\[
x _ {\mathrm{M} / \mathrm{M} ^ {\prime}}. f (m) = \phi (x _ {\mathrm{M}} y _ {\mathrm{M}}. m) = \phi (y _ {\mathrm{M}} x _ {\mathrm{M}}. m) + \phi ([ x, y ] _ {\mathrm{M}}. m).
\]

现在 \([x,y] \in \mathfrak{a}\)，从而 \(\phi([x,y]_{\mathbb{M}}.m) = 0\)；另一方面，

\[
\phi (y _ {\mathrm{M}} x _ {\mathrm{M}}. m) = f (x _ {\mathrm{M}}. m).
\]

因此 \(x_{\mathbf{M} / \mathbf{M}'} . f(m) = f(x_{\mathbf{M}}.m)\)。于是 \(f(\mathbf{M}')\) 是 \(\mathbf{M} / \mathbf{M}'\) 的子 a-模，并同构于 \(\mathbf{M}'\) 的一个商模，因而是 \(\mathbf{N}\) 型纯模；所以 \(f(\mathbf{M}') = \{0\}\)。

推论。设 g 是 K 上的李代数，a 是 g 的理想。设 M 是简单 g-模，并且作为 K-模具有有限长度。则存在简单 a-模 N，使得 M 是 N 型纯 a-模。

由于 \(a\)-模 \(\mathbf{M}\) 具有有限长度，\(\mathbf{N}\) 的子 \(a\)-模集合中存在极小元 \(\mathbf{M}\)：它是 \(a\) 的简单子 \(\mathbf{M}\)-模。因此，\(a\) 中最大的 \(\mathbf{M}\) 型纯子 \(\mathbf{N}\)-模不等于 \(\neq \{0\}\)，并且是 \(g\) 的子 \(\mathbf{M}\)-模（命题 1），所以它必等于 \(\mathbf{M}\)。

